\begin{abstract}
Can points, distances and curvature be built from nothing more than a Markov chain and a decision about what to ignore?
We study this question for finite Markov dynamics within the Six Birds emergence calculus.
A \emph{lens} groups microstates into macro states, each represented by a prototype distribution; lifting, evolving for $\tau$ steps and repackaging gives the macro dynamics.
The distance between two macro states is the cost of the cheapest chain of macro transitions, each step costing the negative logarithm of its symmetrized, floored transition probability.
We prove, and formalize in Lean~4, that this always yields an extended pseudometric, that each macro state's persistence defect equals its escape probability, and that the worst escape probability bounds the closure defect on every microstate.
A layer is \emph{coherent} when its closure, persistence, route and normalized distortion defects all vanish along a declared refinement family.
At stage $\tau=5$ we construct coherent layers, in declared units, for three systems: block lenses on a lazy grid walk, whose normalized metrics converge to the $L^1$ unit square; recursive cells on a Sierpi\'nski gasket walk, converging to a self-similar compact space of Hausdorff dimension $\log 3/\log 2$ in which arbitrarily small balls about every point resist Hilbert approximation to within one sixteenth of their radius; and reservoir chains on the sphere, converging to the round unit sphere, whose curvature a landmark transport estimator recovers from the Markov metrics.
For the isotropic walk on the lattice, centred negative-log probabilities converge uniformly to a quadratic form on diffusive windows, so the square root of the cost, in diffusive units, converges to Euclidean distance and the Pythagorean identity becomes a law of accounting.
The analysis also marks the limits of the approach.
Spectrally learned lenses give valid metrics but fail persistence, for the canonical grid and gasket partitions under every choice of prototype, and a local-MDS holonomy estimator can overestimate curvature by more than $20\%$ in the refinement limit, even on exact spherical distances.
\end{abstract}
